\documentclass{article}
\input{../../lib.sty}

\title{\texttt{fn make\_select\_private\_candidate}}
\author{Michael Shoemate}
\date{}

\begin{document}

\maketitle

\contrib
Proves soundness of \rustdoc{combinators/fn}{make\_select\_private\_candidate} in \texttt{mod.rs}.

\texttt{make\_select\_private\_candidate} constructs a measurement that samples a repetition count,
invokes a private candidate measurement repeatedly, and either returns the first score at least a
threshold or the highest-scoring sampled candidate.

\section{Hoare Triple}
\subsection*{Precondition}
\subsubsection*{Compiler-verified}
\begin{itemize}
    \item Generic \texttt{DI} implements \rustdoc{core/trait}{Domain}.
    \item Generic \texttt{MI} implements \rustdoc{core/trait}{Metric}.
    \item Generic \texttt{MO} implements \texttt{PrivateCandidateMeasure}.
    \item \texttt{(DI, MI)} implements \rustdoc{core/trait}{MetricSpace}.
    \item \texttt{measurement} has type \texttt{Measurement<DI, MI, MO, (f64, TO)>}.
\end{itemize}

\subsubsection*{Caller-verified}
\begin{itemize}
    \item \texttt{measurement} is a valid measurement.
\end{itemize}

\subsection*{Pseudocode}
\lstinputlisting[language=Python,firstline=2,escapechar=|]{./pseudocode/make_select_private_candidate.py}

\subsection*{Postcondition}
\begin{theorem}
\label{postcondition}
\validMeasurement{\texttt{(measurement, mean, threshold, distribution, DI, MI, MO, TO)}}{\texttt{make\_select\_private\_candidate}}
\end{theorem}

\section{Proof}
Assume the preconditions.

\begin{lemma}
\label{construction}
If the constructor returns \texttt{Ok(out)}, then:
\begin{itemize}
    \item \texttt{resolved} satisfies the proof definition of \texttt{ResolvedRepetitions};
    \item the parameter combination is accepted by \texttt{MO::validate};
    \item \texttt{privacy\_map} is the privacy map returned by the corresponding
    \texttt{PrivateCandidateMeasure} implementation.
\end{itemize}
\end{lemma}

\begin{proof}
Lines \ref{line:mean-guard} and \ref{line:threshold-guard} enforce finiteness of the scalar
arguments. Line \ref{line:resolve} invokes
\texttt{Repetitions::resolve}; by its postcondition, a
successful return produces a valid \texttt{resolved}. Line \ref{line:validate} invokes the relevant
\texttt{MO::validate}, so any unsupported combination is rejected before construction succeeds.
Finally, line \ref{line:new-privacy-map} invokes the relevant \texttt{MO::new\_privacy\_map}.
\end{proof}

\begin{lemma}
\label{algorithm}
If the constructor succeeds, the returned function implements the algorithm stated in the proof
definition of \texttt{make\_select\_private\_candidate}.
\end{lemma}

\begin{proof}
By Lemma~\ref{construction}, line \ref{line:sample} samples from a valid
\texttt{ResolvedRepetitions}. By the postcondition of
\texttt{ResolvedRepetitions::sample},
\texttt{remaining} is an exact draw from the selected repetition law.

The loop on line \ref{line:loop} invokes the base \texttt{measurement} exactly once per remaining
repetition. If \texttt{threshold} is present, line \ref{line:threshold-branch} returns
immediately on the first sampled candidate whose score is at least the threshold, and otherwise the
function returns \texttt{None} after the budget is exhausted.

If \texttt{threshold} is absent, line \ref{line:choose} updates \texttt{best}. By the postcondition
of \texttt{choose}, \texttt{best} is always the
highest-scoring non-NaN sampled candidate seen so far, with ties broken in favor of the most recent
candidate. Therefore after the budget is exhausted, the function returns exactly the best sampled
candidate, or \texttt{None} if none was retained.
\end{proof}

\begin{proof}[Proof of data-independent runtime errors.]
Assume the constructor succeeds, and consider any two inputs \texttt{x}, \texttt{x'} in
\texttt{input\_domain}.

Inside the returned function, line \ref{line:sample} samples from \texttt{resolved}. By the
postcondition of \texttt{ResolvedRepetitions::sample}, any runtime error here comes only from the
underlying sampler or numeric casts, and does not depend on \texttt{x}.

The only later fallible operation is line \ref{line:eval}, which invokes the underlying
\texttt{measurement}. By the postcondition of a valid measurement, these runtime errors are
data-independent in \texttt{x}.

The remaining lines are comparisons, calls to \texttt{choose}, and loop control. The proof of
\texttt{choose} shows it is infallible. Therefore the returned function has only data-independent
runtime errors.
\end{proof}

\begin{proof}[Proof of privacy guarantee.]
By Lemma~\ref{construction}, the accepted parameter combination is one supported by the relevant
\texttt{PrivateCandidateMeasure} implementation, and by Lemma~\ref{algorithm} the returned
function is the repeated-selection algorithm claimed by the proof definition.

If \texttt{MO = MaxDivergence}, the postcondition of
\texttt{MaxDivergence} as a \texttt{PrivateCandidateMeasure} shows that the privacy map returned on
line \ref{line:new-privacy-map} is exactly the implemented pure-DP bound for the accepted branch.

If \texttt{MO = RenyiDivergence}, the postcondition of
\texttt{RenyiDivergence} as a \texttt{PrivateCandidateMeasure} shows that the privacy map returned
on line \ref{line:new-privacy-map} is exactly the implemented Rényi-DP bound for the accepted
branch.

The returned measurement reuses the input domain, input metric, and output measure of
\texttt{measurement}, and \texttt{Measurement::new} is invoked with the function and privacy map
just analyzed. Therefore, for every \texttt{d\_in} where the privacy map evaluates to
\texttt{d\_out}, invoking the returned measurement on \texttt{d\_in}-close datasets is
\texttt{d\_out}-close under the output measure.
\end{proof}

\begin{proof}[Proof of Theorem~\ref{postcondition}]
The data-independent runtime error proof and privacy proof establish the two clauses of
\validMeasurement{\texttt{(measurement, mean, threshold, distribution, DI, MI, MO, TO)}}{\texttt{make\_select\_private\_candidate}}.
\end{proof}

\end{document}
